\section{Limitations and Transferability}\label{sec:limitations}
% Derived from thesis chapter 7. Internal infrastructure and capacity details are omitted.
The findings apply to the implemented workflow and tested configurations, with limits on causal attribution and generalisation.

\subsection{Scope of the Comparison}
The small set of differentiated scheduling tools limits transfer to larger tool sets, overlapping capabilities or longer dependencies, where separate planning may be useful.
The missing backend-without-planner variant prevents a complete comparison of both design choices.
Backend orchestration also changes tool granularity, so its effect cannot be attributed to execution location alone.

The selection stage compared three model versions on two standardised scenarios.
With five regular repetitions, one run changes a success rate by 20 percentage points.
Five successes do not establish a reliable population-level success rate.
T01 and T02 supported selection and settings decisions, so their results are not independent confirmation of the selected configuration's performance.
Later functional and robustness cases broadened coverage for Haiku only and cannot establish superiority over the untested alternatives.
Prompts were developed using one evaluated model family, potentially favouring it despite shared prompts and tools across models within each variant.

\subsection{Technical and Evaluation Limits}
Date processing and error handling cover selected expressions and service responses.
Response times also depend on external services.
Latency medians exclude technically failed runs and use different sample sizes, so faster responses must be interpreted alongside success counts.
The limited load tests do not establish sustained production capacity.

Direct prompt-injection tests do not cover indirect attacks or establish comprehensive security.
Ten repetitions of each fixed input assess consistency but do not provide 80 distinct attack patterns.
The functional cases are reported as coverage checks, without repetition counts from which to estimate a success rate.
One case required an implementation adjustment, so the later assessment includes development feedback rather than an unchanged-system validation alone.
Concurrent use does not verify context isolation.
Brief feedback from three operational users cannot establish long-term acceptance, and six developers cannot substitute for operational experience.
Reporting aggregate results without prompts and run-level traces limits independent replication.
Generalised case descriptions, settings and evaluation criteria allow the study design to be inspected, while internal identifiers, deployment details, access mechanisms and operational capacity figures are excluded.

\subsection{Transferability}
Explicit tool boundaries, structured state, deterministic validation, clear results and human confirmation provide design principles for other processes.
Their suitability and the choice of orchestration require task-specific evaluation.

\section{Conclusion and Future Work}\label{Sec:ConclusionAndFutureWork}
% Derived from thesis chapter 8. No additional results are introduced.
This paper presented the design, selection and evaluation of an MCP-based scheduling assistant under human control.
The prototype combines flexible language input with structured dialogue state, bounded tools and application-side validation.
Explicit confirmation keeps the booking decision with the user.

Haiku with individual tools and no separate planner was selected after five correct outcomes per selection scenario.
The subsequent assessment covered seven functional cases, including one that required clearer handling of empty search results.
The 80 direct prompt-injection executions produced no observed control violation, but some disrupted the intended workflow.
Feedback from three operational users indicated that flexible input was useful while concise responses and low delays remained important.
Together, these stages support a reasoned prototype choice and identify development needs, without establishing a universally optimal architecture.
The main lesson is to assess outcome correctness separately from technical execution and to treat configuration selection, functional coverage and user acceptance as distinct sources of evidence.

Future work should complete the comparison with backend orchestration without a planner and introduce independent scenarios, larger tool sets and more repetitions.
Longer user studies, sustained load tests and targeted context-isolation and prompt-injection checks are needed before operational use.
